\subsection{Verifying Completed Rewrites}
\label{sec:correctness}
\label{sec:broader-operators}
Let $C(p)$ be the rule's \code{where} condition and $f$ its attribute program,
so that $\beta=f(p)$. Note that $f$ differs from \code{build} in
Section~\ref{sec:find}: $f$ may be user-written or inferred, whereas
\code{build} is the constructor defined by the mathematical algorithm.
Both are checked under the contract in Equation~\eqref{eq:contract}.

The verifier substitutes the attribute programs and reconstructs the completed
rewrite in a fresh context. It derives coordinate relationships and values
from those programs, rather than assuming conclusions used during synthesis.
Table~\ref{tab:verification-checks} lists the requirements. Only source
legality and the declared guard may be assumed; any additional conditions
needed by construction must follow from them.

\begin{table}[!htbp]
\caption{Checks required for a completed semantic rewrite.}
\label{tab:verification-checks}
\begin{tabularx}{\linewidth}{@{}lY@{}}\toprule
Check & Required property\\\midrule
Attribute computation & Every constructor uses only source metadata, has no
holes, and is evaluable whenever the rule applies.\\
RHS legality & Constructed attributes satisfy operation contracts, including
valid positions, paired extents, reduction domains, and access bounds.\\
Value equivalence & Fully lowered expressions agree in indexed reads,
coefficients, and reduction domains under the selected arithmetic policy.\\
Output agreement & Results have identical ordered output shapes and
compatible fixed dtype.\\\bottomrule
\end{tabularx}
\end{table}

Here $C(p)$ means that evaluating \code{where} successfully returns
\code{true}. Whenever it holds, the attribute program must evaluate without
errors, respecting the language's evaluation order and short-circuiting
rules. Undefined attribute computation cannot satisfy the contract.

Under the theorem's assumptions, using its constructor reduces the
verification obligation to
\begin{equation}
 \forall p,\quad
 \operatorname{Legal}_L(p)\land C(p)\Rightarrow\operatorname{check}(p).
 \label{eq:guard-feasibility}
\end{equation}
Whenever the source is legal and the declared condition holds,
the feasibility check must succeed.
Thus, verification must establish that the declared condition implies
the feasibility check. The theorem then guarantees correctness of
the constructed replacement for every legal match satisfying that condition.
